% Add your abstract here

%Η παρούσα εργασία εστιάζει στην πρακτική μελέτη...\newline \newline

Η μετάβαση από παραδοσιακές, τοπικά φιλοξενούμενες υποδομές σε δυναμικά περιβάλλοντα νέφους αποτελεί επιτακτική ανάγκη για την εξασφάλιση της κλιμακωσιμότητας και της ανθεκτικότητας των σύγχρονων εφαρμογών. Η παρούσα εργασία παρουσιάζει την αρχιτεκτονική εξέλιξη ενός συστήματος Διαχείρισης Έξυπνης Αποθήκης (Smart Warehouse ERP) – το οποίο αρχικά βασιζόταν σε μια στατική, χειροκίνητα αναπτυσσόμενη υποδομή – σε μια πλήρως αυτοματοποιημένη, Cloud-Native πλατφόρμα. Ο κύριος στόχος είναι η επίλυση προβλημάτων περιορισμού πόρων, απώλειας δεδομένων και μη αυτοματοποιημένων κύκλων ανάπτυξης (deployments), μέσω της υιοθέτησης βέλτιστων πρακτικών του κλάδου της Μηχανικής Λογισμικού (DevOps).

Για την επίτευξη αυτού του μετασχηματισμού, υιοθετήθηκε η μεθοδολογία Infrastructure as Code (IaC) μέσω των εργαλείων OpenTofu και Ansible, εξασφαλίζοντας την αυτόματη και επαναλήψιμη εγκαθίδρυση της υποδομής σε περιβάλλον Microsoft Azure. Ο κύκλος ζωής των εφαρμογών (microservices) ενορχηστρώνεται βάσει της φιλοσοφίας GitOps, συνδυάζοντας τον Jenkins για τη Συνεχή Ενσωμάτωση (CI) και το ArgoCD για τη Συνεχή Παράδοση (CD) εντός ενός MicroK8s cluster. Η επιχειρηματική λογική ενισχύθηκε με μοτίβα σχεδιασμού ανθεκτικότητας (Design Patterns), όπως Circuit Breaker και Retry, ενώ ενσωματώθηκε serverless υπολογιστική λογική για την ασύγχρονη, event-driven επεξεργασία δεδομένων. Παράλληλα, εξασφαλίστηκε η μονιμότητα των δεδομένων (Data Persistence) μέσω Managed Disks και επιτεύχθηκε κεντρική διαχείριση ταυτοτήτων και μυστικών μέσω των Keycloak και Vault, αντίστοιχα.

Τέλος, το σύστημα θωρακίστηκε με μηχανισμούς οπτικοποίησης και παρακολούθησης (Observability) με χρήση Prometheus και Grafana. Βάσει πειραματικών δοκιμών φόρτου (Load Testing), εξήχθησαν κανόνες κλιμάκωσης και διαμορφώθηκαν ρεαλιστικές Συμφωνίες Επιπέδου Υπηρεσιών (SLAs), ταξινομημένες σε κατηγορίες Gold, Silver και Bronze. Συνολικά, η εργασία αποδεικνύει εμπράκτως πώς η σύγκλιση των μεθοδολογιών GitOps, IaC και Cloud-Native αρχιτεκτονικών μεταμορφώνει ένα ακαδημαϊκό πρωτότυπο σε μια παραγωγικά έτοιμη (production-grade), πλήρως ανθεκτική πλατφόρμα. \newline \newline

% Keywords
\noindent \textbf{Λέξεις Κλειδιά:} Cloud-Native Architecture, GitOps, Infrastructure as Code (IaC), Continuous Integration and Deployment (CI/CD), Microservices Resilience, Serverless Computing, Service Level Agreements (SLA), Kubernetes Observability.